% BHU PYQ -- Auto-generated & error-corrected
\documentclass[11pt,a4paper]{article}

\usepackage[a4paper,top=2.5cm,bottom=2.5cm,left=2.8cm,right=2.8cm,headheight=14pt]{geometry}
\usepackage{amsmath,amssymb}
\usepackage[T1]{fontenc}
\usepackage[utf8]{inputenc}
\usepackage{lmodern}
\usepackage{microtype}
\usepackage{fancyhdr}
\usepackage{enumitem}
\usepackage{hyperref}
\usepackage{lastpage}
\usepackage{parskip}

\hypersetup{colorlinks=true,urlcolor=black,linkcolor=black,
  pdftitle={BAS-201: Algebra, Probability and Statistical Methods -- BHU PYQ},pdfauthor={Banaras Hindu University}}

\pagestyle{fancy}\fancyhf{}
\renewcommand{\headrulewidth}{0.4pt}\renewcommand{\footrulewidth}{0.4pt}
\fancyhead[L]{\small   \textit{Banaras Hindu University}}
\fancyhead[R]{\small   \textit{BAS-201: Algebra, Probability and Statistica}}
\fancyfoot[C]{\small Page  \thepage\ of \pageref{LastPage}}


\newlist{parts}{enumerate}{2}
\setlist[parts,1]{label=(\alph*),leftmargin=*,itemsep=5pt,topsep=3pt}
\setlist[parts,2]{label=(\roman*),leftmargin=*,itemsep=3pt,topsep=2pt}

\newcommand{\pts}[1]{\hfill   \textit{[\#1]}}


\begin{document}


\begin{center}
  {\large\bfseries BANARAS HINDU UNIVERSITY}\\[2pt]
  {\normalsize B.A. (Hons.) Semester II Examination, 2023-24}\\[6pt]
  \rule{0.8\linewidth}{0.5pt}\\[6pt]
  {\large\bfseries Paper No. BAS-201: Algebra, Probability and Statistical Methods}\\[4pt]
  {\normalsize Time: 3 Hours\hfill Maximum Marks: 70}
\end{center}
\rule{\linewidth}{0.4pt}
\smallskip



\noindent   \textit{Instructions: Answer   \textbf{five} questions including Question~1 which is compulsory. Symbols have their usual meanings.}
\bigskip

Rol] NO. ccsscseesseesseessessseneees
BAS-201 : Algebra, Probability and Statistical Methods
\medskip


\noindent   \textbf{Question 1.}

\begin{parts}
  \item Define the probability of an event. State and prove the addition law of probability for $n$ events. \pts{2+7}
  \item A card is drawn at random from a deck of ordinary playing cards. What is the probability that it is a diamond, a face card, or a king? \pts{5}
\end{parts}
\medskip\rule{\linewidth}{0.2pt}

\medskip


\noindent   \textbf{Question 2.}

\begin{parts}
  \item What do you understand by conditional probability? State and prove Bayes' theorem. \pts{2+7}
  \item If two fair dice are thrown, what is the probability of getting:
    
\begin{parts}
      \item a double six,
      \item a sum of 8 or more dots.
    \end{parts} \pts{5}
\end{parts}
\medskip\rule{\linewidth}{0.2pt}

\medskip


\noindent   \textbf{Question 3.}

\begin{parts}
  \item Explain the concept of a random variable. Differentiate between discrete and continuous random variables. What is a distribution function and what are its properties? \pts{7}
  \item Given the following probability density function:
    $$f(x) = A (4x - 2x^2), \quad 0 < x < 2$$
    and $0$ otherwise. Calculate the value of $A$ so that $f(x)$ is a probability density function. \pts{7}
\end{parts}
\medskip\rule{\linewidth}{0.2pt}

\medskip


\noindent   \textbf{Question 4.}

\begin{parts}
  \item Explain the concepts of Probability Mass Function (p.m.f.) and Probability Density Function (p.d.f.). Define $E(X)$, the expected value of a random variable $X$. \pts{6}
  \item Find the probability distribution of the sum of the dots when two fair dice are thrown. Use the probability distribution to find the probabilities of obtaining:
    
\begin{parts}
      \item a sum of 8 or 11,
      \item a sum that is greater than 8,
      \item a sum that is greater than 5 but less than or equal to 10.
    \end{parts} \pts{8}
\end{parts}
\medskip\rule{\linewidth}{0.2pt}

\medskip


\noindent   \textbf{Question 5.}

\begin{parts}
  \item Define the Binomial distribution and derive its mean and variance. \pts{7}
  \item A fair coin is tossed 5 times. Find the probabilities of obtaining various numbers of heads. \pts{7}
\end{parts}
\medskip\rule{\linewidth}{0.2pt}

\medskip


\noindent   \textbf{Question 6.}

\begin{parts}
  \item Define Poisson distribution and derive its mean and variance. \pts{7}
  \item Write down the steps involved in fitting a Poisson distribution to a given frequency data. \pts{7}
\end{parts}
\medskip\rule{\linewidth}{0.2pt}

\medskip


\noindent   \textbf{Question 7.}

\begin{parts}
  \item Define Normal distribution and derive its mean and variance. \pts{7}
  \item What do you mean by Point Estimation? Explain the desirable properties of a good point estimator. \pts{7}
\end{parts}
\medskip\rule{\linewidth}{0.2pt}

\medskip


\noindent   \textbf{Question 8.}
Explain any   \textbf{four} of the following with examples:

\begin{parts}
  \item Null Hypothesis and Alternative Hypothesis
  \item Simple Hypothesis and Composite Hypothesis
  \item Acceptance Region and Rejection Region
  \item Type-I Error and Type-II Error
  \item One-tailed test and two-tailed test
\end{parts}
\pts{14}

\vfill
\rule{\linewidth}{0.4pt}

\begin{center}\small
  \href{https://bhu-exam.vercel.app/}{Click here to visit our website}\\[4pt]\href{https://me-aryan.vercel.app/}{Click here to visit our contributor}\\[4pt]\href{https://quashan-me.vercel.app/}{Click here to visit our contributor}
\end{center}

\end{document}